% Part: first-order-logic
% Chapter: syntax-and-semantics
% Section: covered-structures

\documentclass[../../../include/open-logic-section]{subfiles}

\begin{document}

\olfileid[hi]{fol}{syn}{cov}

\olsection{प्रथम-क्रम भाषाओं के लिए आच्छादित \printtoken{P}{structure}}

\begin{explain}
स्मरण करें कि यदि किसी पद में कोई !!{variable} न हो तो वह \emph{बंद} है।
\end{explain}

\begin{defn}[बंद पदों का !!^{value}]
यदि $t$, भाषा~$\Lang L$ का कोई बंद पद है और $\Struct M$,~$\Lang L$ की
!!{structure} है, तो \emph{!!{value}}~$\Value{t}{M}$ निम्न प्रकार परिभाषित
है:
\begin{enumerate}
\item यदि $t$ केवल !!{constant}~$c$ है, तो $\Value{c}{M} = \Assign{c}{M}$।
\item यदि $t$ का रूप $\Atom{f}{t_1, \ldots, t_n}$ है, तो
  \[
  \Value{t}{M} = \Assign{f}{M}(\Value{t_1}{M}, \ldots,
  \Value{t_n}{M}).
  \]
\end{enumerate}
\end{defn}

\begin{defn}[आच्छादित !!{structure}]
कोई !!{structure} \emph{आच्छादित} है यदि प्रांत का प्रत्येक अवयव किसी बंद
पद का !!{value} है।
\end{defn}

\begin{ex}
मान लें कि~$\Lang L$ वह भाषा है जिसमें !!{constant} $\Obj{zero}$,
$\Obj{one}$, $\Obj{two}$, \dots, दो-स्थानी !!{predicate}~$<$, और दो-स्थानी
!!{function} $+$ तथा $\times$ हैं। तब !!a{structure}~$\Struct
M$,~$\Lang L$ के लिए, ऐसी लें जिसका प्रांत $\Domain M = \{0, 1, 2, \ldots
\}$ है और नियतन
$\Assign{\Obj{zero}}{M} = 0$, $\Assign{\Obj{one}}{M} = 1$,
$\Assign{\Obj{two}}{M} = 2$, इत्यादि हैं। दो-स्थानी संबंध-चिह्न $<$ के लिए
$\Assign{<}{M}$ उन सभी युग्मों $\tuple{c_1, c_2} \in \Domain{M}^2$ का
समुच्चय है जिनमें $c_1$,~$c_2$ से छोटा है: उदाहरणतः $\tuple{1, 3} \in
\Assign{<}{M}$, किंतु $\tuple{2, 2} \notin \Assign{<}{M}$। दो-स्थानी
!!{function} $+$ के लिए $\Assign{+}{M}$ को सामान्य ढंग से परिभाषित करें---
उदाहरणतः $\Assign{+}{M}(2,3)$ का मान~$5$ है---और दो-स्थानी
!!{function}~$\times$ के लिए भी ऐसा ही करें। अतः $\Obj{four}$ का !!{value}
केवल~$4$ है, और $\times(\Obj{two},
+(\Obj{three},\Obj{zero}))$ का !!{value} (अथवा मध्यस्थ संकेतन में
$\Obj{two} \times
(\Obj{three} + \Obj{zero})$) यह है:
\begin{multline*}
\Value{\times(\Obj{two}, +(\Obj{three},\Obj{zero}))}{M} =\\
\begin{aligned}
& =\Assign{\times}{M}(\Value{\Obj{two}}{M}, \Value{+(\Obj{three}, \Obj{zero})}{M})\\
& = \Assign{\times}{M}(\Value{\Obj{two}}{M}, \Assign{+}{M}(\Value{\Obj{three}}{M},
\Value{\Obj{zero}}{M})) \\
& = \Assign{\times}{M}(\Assign{\Obj{two}}{M}, \Assign{+}{M}(\Assign{\Obj{three}}{M},
\Assign{\Obj{zero}}{M})) \\
& = \Assign{\times}{M}(2, \Assign{+}{M}(3, 0)) \\
& = \Assign{\times}{M}(2, 3) \\
& = 6
\end{aligned}
\end{multline*}
\end{ex}

\begin{prob}
क्या अंकगणित का मानक प्रतिरूप $\Struct N$ आच्छादित है? स्पष्ट कीजिए।
\end{prob}

\end{document}
